\section{Results}

The output files from Serpent were parsed and analyzed using the 
``serpentTools'' python package \cite{johnson_serpenttools_2020}. The 
\keff, neutron flux, and \betaEff were output by Serpent. The reactivity 
feedback coefficients were calculated by fitting a linear function to the 
data (material temperature vs reactivity) and propagating error through 
the fit. The neutron flux was divided into two energy groups for 
fast and thermal fluxes. We applied the default two-group structure in Serpent, 
in which the cutoff between the energy groups is 0.625 eV.

\subsection{Xe-100-like reactor}
\subsubsection{k-eff}

\begin{table}
    \centering  
    \caption{\keff value from each fuel composition in the Sangamon200.}
    \label{tab:xe100_keff}
    \begin{tabular}{c c}
        \hline
        Fuel Composition & \keff \\
        \hline
        Pure & \\
        EBR-II & \\
        Y12 & \\
        \hline        
    \end{tabular}
\end{table}

\subsubsection{Neutron flux}

\subsubsection{Beta-eff}
\begin{table}
    \centering  
    \caption{\betaEff value from each fuel composition in the Sangamon200.}
    \label{tab:xe100_betaEff}
    \begin{tabular}{c c}
        \hline
        Fuel Composition & \keff \\
        \hline
        Pure & \\
        EBR-II & \\
        Y12 & \\
        \hline        
    \end{tabular}
\end{table}

\subsubsection{Reactivity coefficients}
\begin{table}
    \centering  
    \caption{Reactivity feedback coefficients from each fuel composition in the Sangamon200.}
    \label{tab:xe100_feedback}
    \begin{tabular}{c c c c c}
        \hline
        & \multicolumn{4}{c}{Feedback Coefficient} \\
        Fuel Composition & Fuel & Moderator & Coolant & Total \\
        \hline
        Pure & \\
        EBR-II & \\
        Y12 & \\
        \hline        
    \end{tabular}
\end{table}

\subsection{MMR-like reactor}
The results of the MMR-like reactor are presented below. The \keff 
values are presented first, followed by the neutron fluxes, the 
\betaEff, and the reactivity coefficients. 

\subsubsection{k-eff}
The fuel composition has a significant effect on the \keff at each 
burnup step, based on the values presented in Table \ref{tab:mmr_keff}.
Using downblended \gls{EBR} or Y12 fuel results in a higher \keff 
at each burnup step, compared with the pure fuel composition.  

\begin{table}
    \centering  
    \caption{\keff value from each fuel composition in the MMR-like model.}
    \label{tab:mmr_keff}
    \begin{tabular}{c c c c}
        \hline
         & \multicolumn{3}{c}{\keff} \\
        Fuel Composition & Beginning of cycle & Mid-cycle & End of cycle \\
        \hline
        Pure & 1.33803 $\pm$ 0.00022 & 1.18042 $\pm$ 0.00024 & 1.05550 $\pm$ 0.00024\\
        EBR-II & 1.49698 $\pm$ 0.00016 & 1.43844 $\pm$ 0.00018 & 1.39386 $\pm$ 0.00016\\
        Y12 & 1.49666 $\pm$ 0.00017 & 1.43894 $\pm$ 0.00016 & 1.39392 $\pm$ 0.00015\\
        \hline        
    \end{tabular}
\end{table}



\subsubsection{Neutron flux}

Using the downblended fuel compositions results in depressions in the 
thermal (Figure \ref{fig:mmr_thermal_axial}) and fast (Figure 
\ref{fig:mmr_fast_axial}) neutron fluxes through the center of the core. 
We observed a similar depression in the axial flux at the other 
burnup steps. 


\begin{figure}
    \centering 
    \includegraphics[trim = 20 0 0 0,clip,scale=0.8]{thermal_axial_flux_0.pdf}
    \caption{Thermal, axial neutron flux through the MMR-like reactor using 
    each fuel composition at the beginning of life (0 MWd/kgU burnup).}
    \label{fig:mmr_thermal_axial}
\end{figure}

\begin{figure}
    \centering 
    \includegraphics[trim = 20 0 0 0,clip,scale=0.8]{fast_axial_flux_0.pdf}
    \caption{Fast, axial neutron flux through the MMR-like reactor using 
    each fuel composition at the beginning of life (0 MWd/kgU burnup).}
    \label{fig:mmr_fast_axial}
\end{figure}

Additionally, using the downblended fuel compositions resulted in a 
depressed neutron flux (Figures \ref{fig:mmr_thermal_radial} and 
\ref{fig:mmr_fast_radial}) in the radial direction. However, the 
depression in the radial flux is a smaller magnitude than the depression 
in the axial flux. 

\begin{figure}
    \centering 
    \includegraphics[trim = 20 0 0 0,clip,scale=0.8]{thermal_radial_flux_0.pdf}
    \caption{Thermal, radial neutron flux through the MMR-like reactor using 
    each fuel composition at the beginning of life (0 MWd/kgU burnup).}
    \label{fig:mmr_thermal_radial}
\end{figure}

\begin{figure}
    \centering 
    \includegraphics[trim = 20 0 0 0,clip,scale=0.8]{fast_radial_flux_0.pdf}
    \caption{Fast, radial neutron flux through the MMR-like reactor using 
    each fuel composition at the beginning of life (0 MWd/kgU burnup).}
    \label{fig:mmr_fast_radial}
\end{figure}


\subsubsection{Beta-eff}
\begin{table}
    \centering  
    \caption{\betaEff value from each fuel composition in the MMR-like model.}
    \label{tab:mmr_betaEff}
    \begin{tabular}{c c c c}
        \hline
         & \multicolumn{3}{c}{\betaEff} \\
        Fuel Composition & Beginning of cycle & Mid-cycle & End of cycle \\
        \hline
        Pure & 0.00549 $\pm$ 0.00610\\
        EBR-II & 0.00649 $\pm$ 0.00533\\
        Y12 & 0.00644 $\pm$ 0.00475\\
        \hline        
    \end{tabular}
\end{table}


\subsubsection{Reactivity coefficients}

\begin{table}
    \centering  
    \caption{Fuel temperature reactivity feedback coefficients (pcm/K) from 
    each fuel composition in the MMR.}
    \label{tab:mmr_fuel_feedback}
    \begin{tabular}{c c c c}
        \hline
        & \multicolumn{3}{c}{Burnup step} \\
        \hline
        Fuel composition & Beginning of cycle & Mid-cycle & End of cycle \\\hline
        Pure & -2.968 $\pm$ 0.053 & -3.955 $\pm$ 0.144 & -4.613 $\pm$ 0.285\\
        EBR-II & -1.183 $\pm$ 0.041 & -1.191 $\pm$ 0.164 & -1.327 $\pm$ 0.071\\
        Y12 & -1.266 $\pm$ 0.034 & -1.334 $\pm$ 0.026 & -0.989 $\pm$ 0.075\\
        \hline        
    \end{tabular}
\end{table}

\begin{table}
    \centering  
    \caption{Moderator temperature reactivity feedback coefficients from each 
    fuel composition in the MMR.}
    \label{tab:mmr_moderator_feedback}
    \begin{tabular}{c c c c}
        \hline
        & \multicolumn{3}{c}{Burnup step} \\
        \hline
        Fuel composition & Beginning of cycle & Mid-cycle & End of cycle \\\hline
        Pure & -0.871 $\pm$ 0.062 & -1.426 $\pm$ 0.227 & -2.215 $\pm$ 0.258\\
        EBR-II & -0.426 $\pm$ 0.035 & -0.376 $\pm$ 0.173 & -0.336 $\pm$ 0.063\\
        Y12 & -0.354 $\pm$ 0.063 & -0.139 $\pm$ 0.069 & -0.155 $\pm$ 0.100\\
        \hline        
    \end{tabular}
\end{table}

\begin{table}
    \centering  
    \caption{Coolant temperature reactivity feedback coefficients from each fuel composition in the MMR.}
    \label{tab:mmr_coolant_feedback}
    \begin{tabular}{c c c c}
        \hline
        & \multicolumn{3}{c}{Burnup step} \\
        \hline
        Fuel composition & Beginning of cycle & Mid-cycle & End of cycle \\\hline
        Pure & 0.001 $\pm$ 0.067 & -0.085 $\pm$ 0.086 & 0.280 $\pm$ 0.152\\
        EBR-II & 0.083 $\pm$ 0.089 & -0.139 $\pm$ 0.149 & 0.046 $\pm$ 0.076\\
        Y12 & 0.000 $\pm$ 0.048 & 0.001 $\pm$ 0.077 & -0.029 $\pm$ 0.067\\
        \hline        
    \end{tabular}
\end{table}

\begin{table}
    \centering  
    \caption{Total temperature reactivity feedback coefficients from each fuel composition in the MMR.}
    \label{tab:mmr_total_feedback}
    \begin{tabular}{c c c c}
        \hline
        & \multicolumn{3}{c}{Burnup step} \\
        \hline
        Fuel composition & Beginning of cycle & Mid-cycle & End of cycle \\\hline
        Pure & -3.956 $\pm$ 0.106 & -5.473 $\pm$ 0.272 & -6.687 $\pm$ 0.337\\
        EBR-II & -1.554 $\pm$ 0.051 & -1.321 $\pm$ 0.1887 & -1.446 $\pm$ 0.159\\
        Y12 & -1.436 $\pm$ 0.116 & -1.450 $\pm$ 0.204 & -1.521 $\pm$ 0.074\\
        \hline        
    \end{tabular}
\end{table}